\section{Method}
\label{sec:method}

Our update has two parts. Temporal reasoning selects observations and inherited
states that remain valid. The mapper then integrates this selected evidence in
its own representation. We first define the selection and then give its TSDF
realization.

\subsection{Inputs and Temporal State}
\label{sec:problem}

A run receives RGB-D images, camera poses, and object instance masks. A supplied
rigid transform aligns the run with the previous map. Observations of the same
physical object use a consistent identity across runs. We write $\mathcal{O}_k(t)$
for the depth-ray observations received in run $k$ up to time $t$.

Each physical object has a sequence of geometric states, called \emph{fragments}.
A fragment records its reconstructed geometry, source observations, and temporal
interval. Observed motion additionally has a tracked trajectory. A new
reconstruction is kept separately until the system decides whether it refines
the current fragment or belongs to a changed state. Ending a state removes its
geometry from the current view, not from the recorded history.

Let $\mathcal{H}(t)$ contain these state intervals and decisions inferred from
observations available by $t$. It supplies the temporal conditions used below.
The backbone reconstructs surfaces and motion; our update decides which of the
resulting observations and stored states remain eligible for current mapping.

\subsection{Inferring When an Object State Changes}
\label{sec:states}
\label{sec:absence}

We use Khronos for local reconstruction and observed-motion tracking
~\cite{schmid2024khronos}. To test an older state, we project its surface into a
later depth image. Let $z_f$ be its expected depth and $d$ the measured depth at
the same pixel. Measurements close to $z_f$ support the surface. A measurement
behind it is evidence that its old location is empty. A closer measurement
occludes the surface; missing image coverage supplies no evidence of absence.

The object-state test uses 5\,cm surface cells, a 5\,cm depth tolerance, and
views with incidence angles no greater than $60^\circ$. A cell becomes reliable
after at least three identity-consistent observations without a seen-through
contradiction while that state was confirmed present. At each update, the
fraction of tested reliable cells seen through is compared with a Beta presence
model learned from that object's confirmed observations and pooled sensor
statistics. A one-sided log-likelihood ratio against a uniform absence model is
accumulated with a cumulative-sum (CUSUM) test. Absence is declared when its
score exceeds $\ln 99$ and seen-through rays outnumber supporting rays.
% (double check) Keep the finalized presence-model, accumulation, and sampling definitions synchronized with the implementation.

A newly reconstructed object can also establish a different state. Once it has
at least 30 reliable surface cells, we compare it with the inherited fragment.
If most of its surface is within 10\,cm of the inherited surface, it refines
that state; otherwise it starts a new state and ends the previous one. Objects
configured as fixed structure instead accumulate complementary views. These
decisions produce the validity intervals used by geometry integration. For a
hidden transition, the interval boundary is inferred from the before-and-after
observations; the method does not insert an unobserved motion trajectory.

\subsection{Selecting Observations and Memory}
\label{sec:validity}

\paragraph{Observation validity.}
An earlier measurement describes the current map only if the scene state it
measured has not subsequently changed. For an observation $o$ taken at $t_o$,
we define
\begin{equation}
 v_o(x,t)=\mathbf{1}\!\left[\begin{array}{c}
 \text{no inferred state change at }x\\[-1pt]
 \text{during }(t_o,t]
 \end{array}\right].
 \label{eq:temporal_validity}
\end{equation}
This condition is evaluated from $\mathcal{H}(t)$, not from future observations.
Let $b_o(x)$ indicate that $x$ is within the integration support of a valid
measurement ray. Invalid depth, locations outside its image footprint, and
unobserved space beyond its supported range do not contribute. The observation
weight is
\begin{equation}
 \begin{aligned}
 w_o(x,t)&=b_o(x)\,v_o(x,t),\\
 A_k(x,t)&=\sum_{o\in\mathcal{O}_k(t)} w_o(x,t).
 \end{aligned}
 \label{eq:valid_weight}
\end{equation}
Thus $A_k$ counts valid current-run measurements at $x$. If a state begins at
$t_L$, observations from before that state cannot be fused into its current
reconstruction. Observations of an ended state remain associated with its
historical fragment.

\paragraph{When memory remains active.}
Let $I_k(x,t)=1$ where the current object reconstruction identifies $x$ as inside
an observed object body, and $0$ elsewhere. Memory is selected by
\begin{equation}
 m_k(x,t)=\mathbf{1}[A_k(x,t)=0]\,(1-I_k(x,t)).
 \label{eq:memory_gate}
\end{equation}
The first factor gives current valid measurements priority. This includes both
surface measurements and rays that pass through an old surface: observed empty
space is still evidence. The second factor prevents an inherited surface from
remaining inside the body of a currently observed object.
% (double check) Specify the implemented observed-body predicate; do not substitute a bounding-box interior or ground-truth volume for it.

The prior is drawn from inherited states still eligible for the current map.
When object reasoning retires a state, its stored geometry belongs to history
rather than this prior. Without a valid current observation or a conflicting
object-state decision, eligible memory remains unchanged. This avoids turning
lack of observation into evidence of disappearance.

\subsection{TSDF Realization}
\label{sec:tsdf}

For a fixed query time, we omit $t$ to keep the notation compact. Observation
$o$ has camera center $c_o$ and measured range $r_o$ along its ray. Its truncated
signed ray-distance contribution is
\begin{equation}
 \phi_o(x)=\operatorname{clamp}\!\left(r_o-\|x-c_o\|,-T,T\right),
 \label{eq:ray_distance}
\end{equation}
where $T$ is the mapper's truncation distance. The sign is positive before the
measured surface and negative behind it within the supported truncation band.
For depth images stored as optical-axis depth, the reading is first converted
to range along the ray.

Let $D_{k-1}(x)$ and $W_{k-1}(x)$ be the inherited field and its integration
weight, after restricting the prior to eligible states. Applying the temporal
and memory selections gives
\begin{equation}
 D_k(x)=\frac{
 \displaystyle\sum_{o\in\mathcal{O}_k}w_o(x)\phi_o(x)
 +m_k(x)W_{k-1}(x)D_{k-1}(x)}{
 \displaystyle\sum_{o\in\mathcal{O}_k}w_o(x)
 +m_k(x)W_{k-1}(x)}.
 \label{eq:valid_tsdf}
\end{equation}
For the first run, the prior weight is zero. The denominator is the resulting
integration weight $W_k(x)$. If it is zero,
$x$ is unknown and contributes no surface. The current surface is extracted
from the zero level set of $D_k$ where valid support exists.

Equation~\eqref{eq:valid_tsdf} has two direct cases. Where $A_k(x)>0$, the
prior is inactive and the field is determined by current valid observations.
Where $A_k(x)=0$, eligible memory supplies the field unless the object-body
condition excludes it. Old and new evidence are therefore not averaged merely
because they occupy the same spatial cell.

Consider an old chair location. If the current depth ray reaches the wall behind
it, $\phi_o$ is positive at the old surface and $m_k=0$: the inherited chair
surface no longer supplies a zero crossing there. If the location is outside
view or hidden by a nearer surface, it has no such current evidence and eligible
memory remains. The decision is made from the measurements before surface
extraction, including regions where no replacement mesh is produced.

\subsection{Identity and Run Handoff}
\label{sec:contract}

The object-state assignment supplies the identity $\ell_k(x)$ of reconstructed
object geometry; geometry without an object assignment belongs to the background.
Retained memory carries its earlier identity and follows that state's temporal
validity. Object and background geometry use the same evidence-selection rule.
We evaluate the TSDF expression at the fine object resolution and the coarser
online-map resolution. Fine-grid surfaces are used where available, and coarse
surfaces fill regions where the fine grid does not establish a surface. This is
a resolution choice, not a second temporal update rule.
% (double check) Apply identical temporal and memory gates to coarse fill so it cannot reintroduce a retired state.

A run exports the eligible geometric state, integration weights, object-state
history, and statistics required by the temporal inference. The next run loads
this state, aligns new observations to the common frame, and repeats the same
selection and integration. Historical views select their corresponding state
intervals; the current view contains only currently eligible states.

\subsection{Map-Specific Integration}
\label{sec:representation}

The shared part of the method is the validity of observations, the eligibility
of memory, and the temporal identity of map content. Signed-distance averaging
in Eq.~\eqref{eq:valid_tsdf} is the TSDF-specific realization. Another map
representation must expose observation support and execute updates in its own
geometry, while respecting the same temporal selections.

For a point or surfel map, this means retaining or integrating selected surface
measurements. For a Gaussian map, a realization would apply validity to the
observations used for optimization and to the inherited primitives they support.
Retiring geometry must also stop obsolete observations from recreating it during
later optimization. These are the requirements for an adapter, rather than
separate temporal decision rules for each representation.
% (double check) Validate a shared-core Gaussian adapter before claiming cross-representation experimental results.
